\chapter{Le prompt système}
\label{ann:prompt}

Le prompt qui encadre la génération constitue l'un des éléments les plus déterminants du système :
c'est lui qui traduit l'exigence de fiabilité en contraintes que le modèle peut suivre. Sa version
finale impose les règles suivantes.

\begin{enumerate}
  \item Répondre uniquement à partir des articles fournis, sans recourir à aucune autre
        connaissance ni à aucune supposition.
  \item Citer systématiquement les numéros d'articles utilisés, sous la forme
        \texttt{(Article~N)}.
  \item Employer une phrase d'abstention \emph{exacte} lorsque les articles fournis ne permettent
        pas de répondre. La formulation est imposée mot pour mot, afin qu'elle soit détectable
        automatiquement par le banc d'évaluation.
  \item Ne jamais compléter un aveu d'absence par une supposition sur le contenu d'un article
        manquant --- y compris présentée comme approximative ou incertaine.
  \item Répondre \emph{par défaut} lorsque les articles fournis contiennent de quoi le faire,
        l'abstention étant réservée aux cas où le sujet n'est réellement pas traité.
  \item Ne jamais citer un numéro de loi, un numéro d'article, une procédure ou une règle
        provenant d'un autre code, et ne jamais affirmer avoir «~consulté~» un texte : les seuls
        textes disponibles sont les articles fournis.
  \item Ne jamais donner de conseil juridique personnalisé.
  \item Répondre de façon claire et concise, dans la langue de la question.
\end{enumerate}

\section{Genèse des règles 4, 5 et 6}

Ces trois règles n'ont pas été écrites d'emblée : elles sont la trace de défaillances observées, et
leur ordre d'apparition est instructif (tableau~\ref{tab:regles}).

\begin{table}[htbp]
  \centering
  \caption{Ajouts successifs au prompt et leur effet mesuré}
  \label{tab:regles}
  \small
  \begin{tabularx}{\textwidth}{@{}>{\centering\arraybackslash}p{1.1cm} R R@{}}
    \toprule
    \sffamily\bfseries\color{navy}Règle &
    \sffamily\bfseries\color{navy}Défaillance à l'origine &
    \sffamily\bfseries\color{navy}Effet observé \\
    \midrule
    4 & Le modèle constatait l'absence d'un article, puis en devinait le contenu
      & Suppression du comportement, sans effet de bord notable \\
    6 & Hallucination inter-codes : citation d'un «~Code civil marocain~» inexistant
      & Abstention correcte portée à \TauxAbstention~\%, mais refus à tort triplés \\
    5 & Sur-correction introduite par la règle~6
      & Bénéfice de la règle~6 conservé, coût partiellement résorbé \\
    \bottomrule
  \end{tabularx}
\end{table}

La règle~5 est donc un correctif de correctif. Elle illustre un point développé au
chapitre~\ref{chap:evaluation} : durcir un prompt déplace le compromis entre prudence et utilité,
et sans mesure on n'observe que le côté du compromis qu'on regardait déjà.

\section{Phrase d'abstention}

La phrase imposée par la règle~3 est reproduite ci-dessous. Sa formulation exacte importe autant
que son contenu : le banc d'évaluation la reconnaît par appariement, et toute reformulation par le
modèle serait comptée comme une réponse.

\begin{encadre}
«~Je ne dispose pas d'information suffisante dans le corpus fourni (Code du travail, Loi~65-99)
pour répondre à cette question.~»
\end{encadre}

\section{Avertissement joint à chaque réponse}

Toute réponse non abstentionniste est accompagnée de l'avertissement suivant, qui matérialise le
non-objectif énoncé au chapitre~\ref{chap:contexte} : le système restitue de l'information
juridique, il ne conseille pas.

\begin{encadre}
«~Réponse fournie à titre informatif, fondée sur les textes cités ; elle ne constitue pas un
conseil juridique. Pour une situation particulière, consultez un professionnel du droit.~»
\end{encadre}
